\section{Phase 2 --- Calibration from Trading History}
\label{sec:phase2}

\subsection{Why calibrate, and what to calibrate}

Phase~1 was run with heuristic defaults inherited from the original
Avellaneda--Stoikov / Gu\'eant / Cartea--Jaimungal papers: a risk
aversion $\gamma = 0.10$, an arrival-decay $k = 0.30$, a one-session
discount horizon $\omega = 1/(8\!\times\!3600)$, a base quoted size of
70\,000~EUR across ten levels with a geometric decay $\beta = 0.30$,
a 60-second EWMA window for realised volatility, and inventory-tilt
weights $\alpha_{\text{spread}}, \alpha_{\text{imb}}$ picked by
inspection. None of these values are estimated from data: they are
reasonable starting points for the first run but no more. Phase~2
replaces each of them with a value re-estimated from the 30-day
fill history and price record produced by Phase~1.

\paragraph{Scope of the calibration.} Only parameters that govern
\emph{our own quoting behaviour} are re-estimated. Client-side
quantities --- the Poisson arrival intensity $A$, the distance decay
$k_{\text{client}}$ of client demand, the Pareto exponent $\alpha$
and the size range $[x_{\min}, x_{\max}]$ --- are exogenous
properties of the market, not outputs of our strategy. Changing them
between Phase~1 and Phase~2 would mean comparing two different worlds
rather than two different quoters, so they are held identical across
phases and enter the calibration only as fixed inputs to downstream
modules. The parameters that \emph{are} re-estimated belong to the
\texttt{QuoterConfig} object consumed by the market maker:
$\gamma, \omega,$ vol-window, $\alpha_{\text{spread}},
\alpha_{\text{imb}}, k, \beta, Q_{\text{base}},$ and the
staleness threshold stale\_s.

\subsection{Data available from Phase~1}

Every calibrator is fed directly from the artefacts of a completed
Phase~1 run, accessible through the \texttt{Controller} instance:

\begin{itemize}\itemsep2pt
\item \texttt{trade\_history} --- one row per fill, tagged by side,
price, size, ladder level, cross-venue fair mid at inception, post-fill
inventory and whether the leg is a client fill on A or a hedge leg on
B/C.
\item \texttt{step\_log} --- one row per simulation step (\(10^7\)
rows for a 30-day run at 100\,ms) with bid/ask/mid on all three venues,
current inventory, number of resting MM orders, and fills since the
previous step. This is the surface on which quote-distance and
order-flow-imbalance features are reconstructed.
\item \texttt{market\_B.noised\_mid\_price} --- the raw mid-price array
of the deeper reference venue, used directly for realised-volatility
estimation.
\end{itemize}

\subsection{Arrival-intensity decay~$k$}

The A-S optimal-spread formula depends on the sensitivity of the
fill rate to our quoted distance from the mid,
\(\lambda(\delta) = A\,e^{-k\delta}\), where $\delta$ is measured in
basis points. We recover $k$ from the ladder rather than fitting it
to a continuous distance variable: each full reprice posts the ten
levels simultaneously, so the number of observations $n_{\text{posted}}$
is approximately constant across levels and fill counts become a
direct proxy for fill rate. Taking logs gives
\[
\log n_{\text{fills}}(i) \;=\; \log(A \cdot C) \;-\; k \cdot \delta_i,
\]
so an ordinary least-squares regression of log-fill-count on the
mean quoted bps distance at each level yields $-\hat k$ as the slope.
Levels with fewer than three fills are dropped --- too noisy for a
stable estimate --- and the regression is run separately for buys
and sells, with the reported $k$ being their average. The
coefficient is clipped to $[0.01, 5]$ to avoid pathological fits
when fills concentrate at a single level. This approach is a
Cartea--Jaimungal-style empirical recovery of the intensity
parameter from observed fill frequencies, and requires no
likelihood maximisation.

\subsection{Volatility: EWMA span + GARCH(1,1)}

The quoter uses an exponentially-weighted estimator of realised
variance with span \texttt{vol\_window\_s}. Two independent models are
fit and then compared.

\paragraph{EWMA grid search.} We split the 30-day log-return series
on venue~B into a 70\% training segment and a 30\% held-out tail.
For each candidate span $w \in \{500, 1000, 2000, 4000, 6000, 10000\}$
(in simulation steps) we initialise the EWMA variance on the training
tail and then forecast one-step-ahead $\hat h_t$ out-of-sample,
scoring the forecast against the squared return $r_t^2$ via RMSE.
The span that minimises this error is retained. This is equivalent
to picking the RiskMetrics decay that best tracks the realised
variation of the mid --- a standard span-selection procedure that
avoids overfitting the in-sample sample variance.

\paragraph{GARCH(1,1) MLE.} In parallel we fit the classical
Bollerslev specification
\[
h_{t+1} \;=\; \omega_{\text{g}} + \alpha\,r_t^2 + \beta\,h_t,
\]
by maximising the log-likelihood
$\ell = -\tfrac12 \sum_t \left[\log h_t + r_t^2 / h_t\right]$
under the constraints $\omega_{\text{g}} > 0$, $\alpha, \beta \ge 0$
and the stationarity condition $\alpha + \beta < 1$. Optimisation
is performed with a warm start at
$(\alpha_0, \beta_0) = (0.05, 0.94)$, with $\omega_0$ set to the
sample variance times $(1-\alpha_0-\beta_0)$ so that the
unconditional variance matches the data. Because a pure-Python
recursion is slow on $10^7$ observations, the log-return series is
sub-sampled to at most 50\,000 points --- well beyond the sample size
at which GARCH MLE has converged in simulation studies.

\paragraph{Model selection.} The two models are compared on AIC and
BIC. Both are computed under the same conditionally-Gaussian
likelihood: EWMA has one effective parameter (the span), GARCH has
three ($\omega_{\text{g}}, \alpha, \beta$). The smaller-AIC model is
chosen for the quoter. In practice, at a 100\,ms step size EWMA is
selected: the GARCH persistence is high but does not contribute enough
additional log-likelihood to offset the two extra parameters.

\subsection{Spread components $\alpha_{\text{spread}},
\alpha_{\text{imb}}$}

The quoter's reservation-price tilt decomposes into a
$\gamma$-weighted inventory-state widening and an order-flow-imbalance
skew:
\[
\text{half-spread}_t \;=\; \text{A-S}_t
\;+\; \alpha_{\text{spread}}\,\bigl(q_t/K\bigr)^2 \cdot \text{A-S}_t
\;+\; \alpha_{\text{imb}}\,\text{OFI}_t \cdot m_t,
\]
with $q_t$ the current EUR inventory, $K$ the capital and $\text{OFI}_t$
a rolling order-flow imbalance in $[-1, +1]$. To recover the two
$\alpha$ coefficients we work on the effective spread actually earned
at each MM fill, $|p_{\text{fill}} - m_{\text{fair}}|$, and regress it
on the two features evaluated at fill time:
\[
|p_{\text{fill}} - m_{\text{fair}}|
\;=\; \beta_0
\;+\; \beta_1 \bigl(q_t/K\bigr)^2
\;+\; \beta_2 \,\text{OFI}_t
\;+\; \varepsilon.
\]
OFI is reconstructed from the fill stream itself via a rolling
fifty-fill window, as twice the fraction of buys in the window minus
one, so a streak of client buys pushes OFI towards $+1$. The
calibrated coefficients are mapped back to the quoter's
parametrisation through $\hat\alpha_{\text{spread}} = \beta_1 / \beta_0$
and $\hat\alpha_{\text{imb}} = |\beta_2| / \bar m_{\text{fair}}$,
which divides out the units and puts both quantities on the same
scale as the heuristic defaults. Both values are then clipped to
plausible ranges ($[0, 5]$ and $[0, 0.01]$ respectively) as a
safeguard against OLS artefacts on the short side of the sample.

\subsection{Staleness threshold stale\_s}

A resting quote should be pulled when the probability of seeing any
further fill before the next price move has become negligible. The
best level is the busiest on the ladder, so if fills at level~1
arrive as a Poisson process with intensity $\lambda_1$, the time to
next fill is exponential with mean $1/\lambda_1$ and
$\mathbb{P}(T > t) = e^{-\lambda_1 t}$. Setting this probability to
$25\%$ gives
\[
t_{75} \;=\; \frac{\log 4}{\lambda_1},
\]
which we use directly as stale\_s. The rate $\lambda_1$ is estimated
simply as the total number of level-1 fills divided by the total
elapsed simulation time. Deeper levels have lower arrival rates, so
level~1 yields the tightest (and most conservative) threshold ---
deep quotes that would technically still be alive at $t_{75}$ are
handled by the next full reprice anyway. The result is clipped to
$[0.5, 60]$~seconds to avoid outliers from degenerate samples.

\subsection{Risk aversion $\gamma$ and horizon $\omega$}

This is the only calibration step that combines the outputs of the
others. Rather than trying to write a likelihood for our own
decisions (the quoter is a deterministic policy, not a probabilistic
model), we pick $(\gamma, \omega)$ by maximising a closed-form
estimate of steady-state utility per second:
\[
U(\gamma, \omega) \;=\;
\underbrace{\lambda(\delta^\ast) \cdot \delta^\ast \cdot \mathbb{E}[X]}_{\text{spread P\&L}}
\;-\;
\underbrace{\gamma\,\sigma^2
\cdot \frac{1}{\omega\,T_{\text{year}}}
\cdot \lambda(\delta^\ast)\,\mathbb{E}[X^2]\,\frac{1}{T_{\text{year}}}}_{\text{inventory-variance cost}},
\]
where $\delta^\ast(\gamma,\omega)$ is the A-S optimal half-spread
replicated from the quoter,
\[
\text{spread}^{\text{AS}}
\;=\; \gamma\,\sigma^2\,\frac{1}{\omega\,T_{\text{year}}}
\;+\; \frac{2}{\gamma}\log\!\bigl(1 + \gamma / k\bigr).
\]
The ingredients are taken from the earlier calibrators and from the
\texttt{ClientFlowConfig} defaults: $\sigma$ is the annualised log-return
standard deviation computed on the venue-B mid; $(A_{\text{buy}}, k_{\text{buy}},
A_{\text{sell}}, k_{\text{sell}})$ are the exogenous client arrival
parameters; $\mathbb{E}[X]$ is the mean of the truncated Pareto client
size distribution, evaluated in closed form from
$(\alpha, x_{\min}, x_{\max})$.

The search proceeds in two stages. We first evaluate $U$ on a
coarse grid of seven candidate values for $\gamma$ (from $0.01$ to
$2.0$) crossed with three candidate horizons for $\omega$
(corresponding to 24-hour, 8-hour and 4-hour windows), and keep the
pair $(\gamma^\ast, \omega^\ast)$ that maximises utility. We then
refine $\gamma^\ast$ by a one-dimensional search over a narrow
bracket around it, holding $\omega^\ast$ fixed, which sharpens the
estimate without straying far from the best grid point. The whole
procedure is analytical on the historical path --- it does not
re-run the simulator --- so it is cheap enough to sit inside the
calibration pipeline. An expected daily Sharpe ratio is reported at
the optimum for diagnostic purposes.

\subsection{Ladder shape $Q_{\text{base}}$ and $\beta$}

The ladder parameters do not lend themselves to an MLE fit: the
simulator never lets us observe what would have happened had we
quoted with a different $Q_{\text{base}}$, and $\beta$ controls
size decay across levels that Phase~1 barely populates. We therefore
derive $Q_{\text{base}}$ from a design input rather than a data fit.
The ladder-level sweep-through rate is the probability that a client
order is large enough to consume the entire first level:
\[
\mathbb{P}(X > Q_{\text{base}})
\;=\; \left(x_{\min}/Q_{\text{base}}\right)^{\alpha}
\;=\; \text{target\_sweep},
\]
under the truncated Pareto with shape $\alpha$ and floor $x_{\min}$.
Inverting this yields
$Q_{\text{base}} = x_{\min} \cdot \text{target\_sweep}^{-1/\alpha}$.
With $\alpha = 1.5$, $x_{\min} = 1\,000$~EUR, and a target sweep
rate of 5\%, the calibrated $Q_{\text{base}}$ drops from 70\,000~EUR
to roughly 7\,370~EUR. The shape parameter $\beta$ is chosen by hand
(default $\beta = 1$ in Phase~2, so level~2 quotes $\approx 37\%$ of
level~1). The choice reflects a prior on ladder concentration rather
than an estimated quantity.

\subsection{How the calibrators fit together}

The five estimators and the ladder derivation are run as a single
pipeline that reads Phase~1 data once and produces the full set of
Phase~2 quoter parameters. The dependency graph is shallow. Four
estimators --- volatility, spread components, arrival decay $k$ and
staleness threshold --- can be computed independently of each other,
because each relies on a different slice of the Phase~1 output
(mid-price series, effective-spread regression, fills by level,
level-1 fill rate). A second stage then uses the volatility output
to compute $\sigma$ for the $\gamma$-optimisation, and inverts the
Pareto client-size distribution to set the base quoted size
$Q_{\text{base}}$.

\begin{figure}[ht]
\centering
\begin{tabular}{rcl}
\multicolumn{3}{c}{\textbf{Stage 1 --- independent estimators}} \\[3pt]
Volatility           & $\longrightarrow$ & $\sigma,$ vol-window \\
Spread components    & $\longrightarrow$ & $\alpha_{\text{spread}}, \alpha_{\text{imb}}$ \\
Arrival decay        & $\longrightarrow$ & $k$ \\
Staleness threshold  & $\longrightarrow$ & stale\_s \\[5pt]
\multicolumn{3}{c}{\textbf{Stage 2 --- depends on Stage 1}} \\[3pt]
Risk-aversion optimisation ($\sigma,$ arrival, size)  & $\longrightarrow$ & $\gamma, \omega$ \\
Pareto-inversion ($\alpha, x_{\min},$ target-sweep)   & $\longrightarrow$ & $Q_{\text{base}}$ \\
\end{tabular}
\caption{Dependency structure of the Phase~2 calibration pipeline.
Stage~1 estimators are mutually independent; Stage~2 consumes the
volatility output together with the fixed client-flow parameters.}
\label{fig:p2-deps}
\end{figure}

Once the pipeline finishes, the calibrated values replace the
Phase~1 heuristics in the quoter configuration, and the Phase~2
backtest is launched with exactly the same simulation machinery as
Phase~1 --- fresh stock and market paths, identical client-flow
parameters --- with only the quoting configuration having changed.

\subsection{Calibrated parameter values}

Running the pipeline on the 30-day Phase~1 run produces the values
shown in Table~\ref{tab:p2-calibrated}. The movements relative to the
heuristic defaults are informative in their own right.

\begin{table}[ht]
\centering
\begin{tabular}{lrr}
\toprule
Parameter & Heuristic default & Calibrated \\
\midrule
$\gamma$                      & $0.10$        & $0.0080$ \\
$\omega$                      & $3.47\!\times\!10^{-5}$ & $1.16\!\times\!10^{-5}$ \\
vol\_window\_s                & $60$          & $1.25$ \\
$\alpha_{\text{spread}}$      & $0.5$         & $5.0$ \\
$\alpha_{\text{imb}}$         & $2.0\!\times\!10^{-4}$ & $1.64\!\times\!10^{-5}$ \\
$k$                           & $0.30$        & $0.30$ \\
$\beta$                       & $0.30$        & $1.00$ \\
$Q_{\text{base}}$ (EUR)       & $70{,}000$    & $7{,}368$ \\
stale\_s                      & $10$          & $60$ \\
\bottomrule
\end{tabular}
\caption{Heuristic vs calibrated \texttt{QuoterConfig} parameters.
EWMA is selected over GARCH(1,1) on AIC. The expected daily Sharpe
computed at the utility optimum is $508$, which is a diagnostic upper
bound on a frictionless path rather than a realised backtest figure.}
\label{tab:p2-calibrated}
\end{table}

Three movements dominate. First, $\gamma$ drops by more than an order
of magnitude: the heuristic penalty was far too aggressive given the
realised mid-price volatility on venue~B, and the optimiser flattens
the inventory-variance term until expected spread capture dominates.
Second, $Q_{\text{base}}$ collapses from 70\,000~EUR to roughly
7\,400~EUR because the Pareto inversion with a 5\% sweep target
simply does not support the original level-1 size. Third,
$\alpha_{\text{spread}}$ climbs from $0.5$ to the clipping boundary
at $5.0$: the regression detects a large inventory-widening effect
in the effective-spread distribution of Phase~1 fills, which the
hand-picked default had under-weighted. The arrival decay $k$ is
unchanged at $0.30$ because the buy/sell log-regressions recover
essentially the same slope the heuristic assumed --- a welcome
confirmation that the shape of the client-demand curve was already
well-specified. EWMA is preferred over GARCH(1,1) on AIC, consistent
with the earlier remark that at 100\,ms step size the GARCH
persistence does not earn back its two extra parameters.

\subsection{Backtest --- Phase~1 vs Phase~2}

The calibrated configuration is evaluated in two ways. First, as a
\emph{controlled} sanity check, it is replayed on the exact 30-day
Phase~1 stock path and client seed --- only the \texttt{QuoterConfig}
changes --- so any delta isolates the strategy improvement from
market randomness. Second, it is rerun on a fresh 15-day stock path
(different seed, finer $dt = 50$\,ms) to confirm that the gains are
not an artefact of the single calibration path.
Table~\ref{tab:p1-vs-p2} reports the headline decomposition for both.

\begin{table}[ht]
\centering
\begin{tabular}{lrrr}
\toprule
Metric & Phase~1 (30d) & Phase~2 same path & Phase~2 fresh 15d \\
\midrule
Total MtM P\&L (USD)         & $+7{,}464$    & $+8{,}935$    & $+5{,}761$ \\
Realised cash P\&L (USD)     & $-139{,}817$  & $-308{,}256$  & $+61{,}629$ \\
Unrealised P\&L (USD)        & $+147{,}281$  & $+317{,}191$  & $-55{,}869$ \\
Inception-spread P\&L (USD)  & $+8{,}415$    & $+10{,}224$   & $+5{,}569$ \\
Total fees (USD)             & $4{,}575$     & $3{,}087$     & $1{,}808$ \\
MM fills                     & $14{,}375$    & $27{,}201$    & $14{,}084$ \\
Hedge legs                   & $10$          & $8$           & $6$ \\
Final inventory (EUR)        & $+136{,}838$  & $+294{,}701$  & $-49{,}104$ \\
\bottomrule
\end{tabular}
\caption{Phase~1 (heuristic) vs Phase~2 (calibrated) on the same
30-day stock path, and Phase~2 replayed on a fresh 15-day path
($dt = 50$\,ms) with identical client-flow parameters.}
\label{tab:p1-vs-p2}
\end{table}

Three observations stand out. First, on the same path the calibrated
quoter posts roughly 90\% more fills (27\,201 vs 14\,375) because
$Q_{\text{base}}$ drops from 70\,000 to 7\,368~EUR, fragmenting each
ladder level into a size that matches the tail of the Pareto client
size distribution. Paradoxically, \emph{total fees fall} from
4\,575 to 3\,087~USD despite the extra fill count: the finer ladder
captures smaller notional per fill, so the maker-fee base shrinks
even as the number of print events grows. Second, inception-spread
P\&L --- the quantity that isolates the quality of our quoting net
of path-dependent drift --- rises from 8\,415 to 10\,224~USD on the
same path, a clean $+21\%$ improvement attributable purely to the
calibrated $\alpha_{\text{spread}}$, $\alpha_{\text{imb}}$ and lower
$\gamma$. Total MtM P\&L rises in tandem (7\,464 $\to$ 8\,935~USD,
$+19.7\%$). Third, the fresh 15-day rerun corroborates the
direction: the calibrated quoter earns $+5\,761$~USD MtM on a
different stochastic path with comparable inception-spread density
per day ($\approx 371$~USD/day vs $\approx 341$~USD/day on the
Phase~1 path), hedges fire six times on a path whose volatility
breached the 90\% delta limit, and the final inventory of
$-49\,104$~EUR is well below a third of the raw Phase~1 figure --- a
product of the refined OFI skew rather than more aggressive risk
offload. The calibration achieves what it was designed for: more
fills at better prices, with inventory kept closer to neutral and
total MtM P\&L materially higher on any path we tested.

\subsection{Diagnostics on the calibrated run}

Two diagnostic plots exported from the calibration notebook give
intuition for \emph{how} the Phase~2 quoter uses its parameters
across the 15-day fresh path.

\begin{figure}[ht]
\centering
\includegraphics[width=0.95\linewidth]{phase2_2.png}
\caption{Quoted spread on A (mean $8.04$~bps) against the
reference-venue spread on B (mean $2.48$~bps) and the annualised
realised volatility of the B mid. The A spread widens visibly during
the three vol spikes on the path and compresses back toward its
floor in the quiet segments, while staying everywhere above the B
reference spread.}
\label{fig:p2-spread-vs-vol}
\end{figure}

Figure~\ref{fig:p2-spread-vs-vol} is the cleanest visual check that
the calibration is doing its job. The quoted spread on A tracks the
realised-vol envelope on B --- precisely the behaviour the A-S
$\gamma\sigma^2/\omega$ term plus the calibrated EWMA window are
designed to produce --- and it never collapses below the reference
spread on B, as a well-behaved DMM quote should not. The
persistence of a $\sim$5--6~bps wedge between the two curves in calm
conditions is the inventory-risk premium the quoter charges on top
of the pass-through of B's spread, not a mis-calibration.

\begin{figure}[ht]
\centering
\includegraphics[width=0.95\linewidth]{phase2_3.png}
\caption{Fill rate and order flow on the Phase~2 fresh-path run.
Top-left: rolling 10-min fill rate against the theoretical Poisson
intensity $A$. Top-right: empirical fill-size distribution,
consistent with the truncated-Pareto client model. Bottom-left:
50-fill rolling order-flow imbalance mean-reverting around zero.
Bottom-right: fill share per ladder level --- level~1 dominates,
as expected given the ladder-size calibration.}
\label{fig:p2-fill-flow}
\end{figure}

Figure~\ref{fig:p2-fill-flow} validates the client-side assumptions
the calibration rests on. The rolling fill rate hugs the
theoretical Poisson intensity $A$, so the arrival process we
regress $k$ against is well-specified; the fill-size histogram
inherits the heavy right tail of the truncated-Pareto input, which
is what the $Q_{\text{base}}$ inversion presumes. The OFI series
mean-reverts around zero over 50-fill windows --- necessary for
$\alpha_{\text{imb}}$ to act as a skew rather than a persistent
drift term --- and level~1 dominates the fill share, confirming
that the ladder-concentration implied by $\beta$ matches where the
realised flow actually lands.